
Para experimentar respecto al uso de hardware de tiempo real y la ejecución del modelo de Hindmarsh-Rose, se ha creado un programa llamado \textit{Pico Neuron} (Repositorio \ref{SEC:PICONEURONREPO}). Dicho programa crea un archivo binario que puede ser introducido en una Raspberry Pi Pico W o en una Raspberry Pi Pico 2 (dependiendo del parámetro que se le pase al compilar). Dicho programa ejecuta un modelo neuronal y proporcional el voltaje generado por el mismo utilizando 6 pines (3 para una salida y 3 para otra). El único problema es que esta salida está limitada por la tasa de baudios, pero es una herramienta que merece la pena explorar para el diseño de circuitos biohíbridos debido a su bajo coste y su amplia flexibilidad.

Dada la naturaleza de la Raspberry Pi Pico, todos los parámetros de configuración van dentro de un archivo, definiendo las variables en tiempo de compilación. La elección de archivo se realiza pasando un parámetro al compilar. Además, en el mismo apéndice que contiene el repositorio de \textit{Pico Neuron}, se encuentra \textit{Pico Neuron Serial} (Repositorio \ref{SEC:PICONEURONREPO}), un programa que lee un puerto con un hilo en tiempo real para obtener los datos generados por el modelo de la Raspberry Pi Pico (Figuras \ref{FIG:HRREG}, \ref{FIG:HRCHAO}).